\documentclass[11pt, a4paper]{article}

\usepackage[T1]{fontenc}
\usepackage[utf8]{inputenc}
\usepackage{tgpagella}
\usepackage[spanish, es-noshorthands]{babel}
\usepackage{amsmath, amssymb, bm}
\usepackage{booktabs}
\usepackage{tabularx}
\usepackage[table, xcdraw]{xcolor}
\usepackage[most]{tcolorbox}
\usepackage[american]{circuitikz}
\usepackage[margin=2.5cm]{geometry}
\usepackage{fancyhdr}

\pagestyle{fancy}
\fancyhf{}
\setlength{\headheight}{16pt}
\lhead{\textbf{UCB Sede Tarija} | Ing. Mecatrónica}
\chead{}
\rhead{IMT-121 Circuitos Electrónicos I | Ejercicios S06-3}
\lfoot{Prof: M.Sc. Ing. Bernardo Quiroga Turdera}
\rfoot{Página \thepage}
\renewcommand{\headrulewidth}{0.4pt}

\definecolor{ucbblue}{RGB}{0, 51, 102}

\begin{document}

\begin{tcolorbox}[colback=ucbblue!5!white, colframe=ucbblue, title=\textbf{Práctica Independiente: Conceptos y Teoremas Terminales}]
\end{tcolorbox}

\section*{Ejercicio 1: Identificación Conceptual de Puerto}
Para un equipo de audio (carga $R_L$) conectado a la salida de un amplificador de varias etapas (terminales $a$ y $b$):
\begin{enumerate}
    \item Trace un diagrama de bloques sencillo marcando: Red Interna, Carga, Puerto, $V_{ab}$ e $I_L$.
    \item Indique exactamente dónde colocaría un multímetro para medir experimentalmente $V_{th}$.
    \item Indique dónde se colocaría un puente (cable) para determinar $I_N$.
\end{enumerate}

\vspace{0.5cm}

\section*{Ejercicio 2: Equivalencia Analítica}
Dada la siguiente red interna conectada a un puerto $a-b$:
\begin{itemize}
    \item $V_s = 9\,\text{V}$ (polo positivo arriba, negativo a tierra).
    \item $R_1 = 3\,\text{k}\Omega$ (entre el terminal positivo de $V_s$ y el nodo $a$).
    \item $R_2 = 6\,\text{k}\Omega$ (entre el nodo $a$ y tierra).
    \item $I_s = 0.5\,\text{mA}$ (inyectada desde la tierra hacia el nodo $a$).
    \item El terminal $b$ está conectado directamente a tierra.
\end{itemize}

\textbf{Tareas:}
\begin{enumerate}
    \item \textbf{Hallar $\bm{V_{th}}$:} Dibuje el circuito con el puerto $a-b$ en circuito abierto y encuentre la tensión del nodo $a$.
    \item \textbf{Hallar $\bm{R_{th}}$:} Redibuja el circuito desactivando adecuadamente ambas fuentes y calcula la resistencia equivalente vista desde el puerto $a-b$.
    \item \textbf{Hallar $\bm{I_N}$:} Usa la relación matemática de conversión para encontrar la fuente de corriente de Norton.
    \item \textbf{Esquemas:} Dibuja el Circuito Equivalente de Thévenin y el Circuito Equivalente de Norton de esta red.
    \item \textbf{Verificación Operativa:} Conecta imaginariamente una carga externa $R_L = 6\,\text{k}\Omega$ a tu modelo de Thévenin. Calcula la tensión final $V_L$ que aparecería en esa carga.
\end{enumerate}

\end{document}
